% ============================================================
% sec/04metodologia.tex
%
% DEBE CONTENER:
%   - Formalización del problema (formulación matemática si aplica)
%   - Pipeline previsto
%   - Protocolo experimental: particiones, validación, métricas
%     Y SU JUSTIFICACIÓN
%   - Plan de baselines
%   - Riesgos identificados con su mitigación
%
% Cada decisión sin justificación técnica cuesta -5 puntos,
% hasta un máximo de -15.
% ============================================================
\section{Metodología Propuesta}
\label{sec:metodologia}

\subsection{Formalización del problema}
\label{subsec:formalizacion}
\pc{Tipo de problema. Definir el espacio de entrada, el espacio de
salida, la función objetivo y las restricciones.}

\begin{equation}
\label{eq:objetivo}
\pc{formulación}
\end{equation}

\noindent\textbf{Definición del target.} \pc{Cómo se deriva la
etiqueta, y qué casos se excluyen y por qué.}

\noindent\textbf{Unidad de observación.} \pc{Qué es una fila.}

\noindent\textbf{Unidad independiente.} \pc{Cuál es la unidad
estadísticamente independiente. Si difiere de la unidad de observación,
ESA distinción gobierna todas las particiones e intervalos de confianza
de este trabajo.}

\subsection{Pipeline previsto}
\label{subsec:pipeline}
\pc{Descripción de extremo a extremo: ingesta, preprocesamiento,
construcción de características, modelo, evaluación, y la interfaz que
consumirá el agente de la Etapa 3.}

\begin{figure}[!t]
  \centering
  \safeimage[width=\columnwidth]{pipeline.png}
  \caption{Pipeline propuesto. \pc{leyenda}}
  \label{fig:pipeline}
\end{figure}

\subsection{Protocolo experimental}
\label{subsec:protocolo}

\subsubsection{Particiones}
\pc{Proporciones exactas. Declarar si la partición es aleatoria,
estratificada, temporal o agrupada, y JUSTIFICARLA frente a las
alternativas descartadas.}

\subsubsection{Estrategia de validación}
\pc{Holdout, k-fold, validación cruzada agrupada o anidada. Declarar
cómo se seleccionan los hiperparámetros sin tocar el conjunto de
prueba, y cuántas veces se toca ese conjunto.}

\subsubsection{Métricas y su justificación}
\begin{table}[!t]
\centering
\caption{Métricas de evaluación y su justificación.}
\label{tab:metricas}
\footnotesize
\begin{tabularx}{\columnwidth}{@{}P{2.2cm} Y@{}}
\toprule
\textbf{Métrica} & \textbf{Por qué esta y qué alternativa desplaza}\\
\midrule
\pc{principal}  & \pc{justificación}\\
\pc{secundaria} & \pc{justificación}\\
\pc{secundaria} & \pc{justificación}\\
\bottomrule
\end{tabularx}
\end{table}

\pc{Un párrafo explicando por qué la métrica principal es la principal,
y qué queda sin capturar que obliga a reportar las secundarias.}

\subsection{Plan de baselines}
\label{subsec:planbaselines}

\pc{Cada nivel se justifica por lo que el anterior no logró explicar,
no por moda. La ausencia de baselines cuesta -10 puntos aunque el
modelo principal funcione.}

\begin{table}[!t]
\centering
\caption{Escalera de baselines.}
\label{tab:baselines}
\footnotesize
\begin{tabularx}{\columnwidth}{@{}c P{2.4cm} Y@{}}
\toprule
\textbf{Nivel} & \textbf{Modelo} & \textbf{Pregunta que responde}\\
\midrule
B0 & \pc{trivial}       & \pc{¿cuál es el piso?}\\
B1 & \pc{heurística}    & \pc{¿hace falta ML?}\\
B2 & \pc{lineal}        & \pc{¿la relación es lineal?}\\
B3 & \pc{no lineal}     & \pc{¿hay interacciones?}\\
B4 & \pc{ensamble}      & \pc{¿cuánto compra la complejidad?}\\
\bottomrule
\end{tabularx}
\end{table}

\noindent\textbf{Tercera referencia.} El baseline de literatura
declarado en la Sección~\ref{sec:baseline} se arrastra como referencia
externa junto a los baselines internos anteriores.

\subsection{Contrato de reproducibilidad}
\label{subsec:reproducibilidad}
\pc{Semillas fijas y dónde viven, configuraciones guardadas, registro
de ejecuciones, hardware, dependencias declaradas, y el comando exacto
que reproduce cada número reportado. Un tercero tiene que poder
reproducir los resultados sin preguntarnos nada.}

\subsection{Riesgos y mitigación}
\label{subsec:riesgos}
\begin{table}[!t]
\centering
\caption{Riesgos identificados y su mitigación.}
\label{tab:riesgos}
\footnotesize
\begin{tabularx}{\columnwidth}{@{}P{2.0cm} c Y@{}}
\toprule
\textbf{Riesgo} & \textbf{Impacto} & \textbf{Mitigación / disparador}\\
\midrule
\pc{riesgo} & \pc{A/M/B} & \pc{mitigación}\\
\pc{riesgo} & \pc{A/M/B} & \pc{mitigación}\\
\pc{riesgo} & \pc{A/M/B} & \pc{mitigación}\\
\pc{riesgo} & \pc{A/M/B} & \pc{mitigación}\\
\bottomrule
\end{tabularx}
\end{table}